\begin{textsample}{POS Dim 2 – ap – Score 28.00 – ap/ap\_inf\_e\_ef\_9.txt}  \label{ex:f2_pos_087}
2.5.
Modalidades de Ensino – Educação de Jovens e Adultos “ Educação de Jovens e Adultos, como direito deve estar articulado com a cidadania, o respeito à diversidade étnico racial, cultural, de gênero, social, ambiental e regional ; com perspectiva interdisciplinar e \textbf{intersetorial} ” ( PAIVA 2014 ).
A Educação Jovens e Adultos - EJA é uma modalidade primogênita no sistema educacional, resultante das lutas sociais em prol da classe \textbf{trabalhadora}, caracterizada por uma proposta pedagógica \textbf{flexível}, não direcionada apenas à especificidade etária, mas, primordialmente, à questão cultural, às diferenças individuais e aos conhecimentos informais adquiridos pelos alunos a partir das suas experiências de vida.
A EJA fundamenta -se na LDBEN 9.394/96, em seu \textbf{art}. 37, como modalidade regular da educação básica, avalizada como direito ao longo da vida[^5 ], pela V Conferência Internacional de Educação de Adultos-CONFINTEA ( 1997 ), no sentido de desenvolver a autonomia e a promoção da igualdade entre homens e mulheres, considerando suas características e especificidades, interesses, condições de vida e perspectivas de trabalho, mediante \textbf{cursos} e \textbf{exames}, viabilizando o acesso e a permanência do \textbf{trabalhador} na escola, com ações integradas e complementares ( BRASIL, 1996 ; UNESCO, 1997 ).
Diante do exposto, é preciso romper com a padronização do ensino regular e reconhecer o contexto e motivações dos alunos da EJA, que na sua maioria retorna aos estudos por uma necessidade, na perspectiva de melhorar sua condição de vida, seja pela inserção no mercado de trabalho ou pela ascensão funcional.
Esse rompimento deve realizar -se na maneira de conduzir o processo educativo e a forma de avaliar os sujeitos da EJA, que abrange da mulher dona de casa ao \textbf{trabalhador}, dos analfabetos absolutos aos analfabetos funcionais, das populações indígenas, do campo, remanescentes de quilombo, pessoas com necessidades específicas e privadas de liberdade.
Esta nova concepção apresenta desafios às práticas existentes, requerendo maior relacionamento entre os sistemas formais e não formais ; além de inovação, criatividade e flexibilidade, a partir da compreensão do universo cultural dos alunos.
O objetivo principal é criar uma sociedade instruída, comprometida com a justiça e o bem estar social, discutida numa perspectiva dialógica com os oprimidos e buscando o rompimento com a opressão, no desejo de construir uma sociedade com equidade ( UNESCO, 1997 ).
Neste contexto, o \textbf{Parecer} 11/2000 do Conselho Nacional de Educação - CNE, refere -se à Educação de Jovens e Adultos, como dívida social não reparada, o que gera uma nova função, entendida como “ direito promotor da cidadania ”, as demais funções garantem os direitos sociais de acesso e permanência ( Equalizadora ), a educação ao longo da vida ( Reparadora ) e a necessidade constante de atualização e de aprendizagem ( Qualificadora ) ( BRASIL, 2000 ).
A Educação de Jovens e Adultos - EJA poderá ser ministrada em estabelecimentos de ensino da rede pública e privada, em \textbf{conformidade} com o \textbf{disposto} na Lei 9394/96 ; nas \textbf{Diretrizes} Curriculares Nacionais para o Ensino Fundamental e Médio ; nas Resoluções CNE/CEB no 01/2000, 03/2010, 04/2010, 07/2010, 02/2012 e 08/2012 e no \textbf{disposto} na Lei no 8069/90, do Estatuto da Criança e do Adolescente ( BRASIL, 1990, 2000, 2010a, 2010b, 2010c, 2012a e 2012b ).
Os \textbf{cursos} da EJA terão estrutura e metodologias específicas, como \textbf{calendário} escolar, turnos e \textbf{horários}, considerando seus objetivos e as características dos educandos, otimizando o tempo escolar e garantindo o processo educativo.
A \textbf{oferta} da EJA pode ser presencial ou à distância, quando credenciados pelo Conselho Estadual de Educação - CEE/AP, com exceção das 1ª e 2ª etapas ( AMAPÁ, 2015a ).
Em se tratando da \textbf{oferta} de educação para as pessoas privadas de liberdade, está expressa o reconhecimento do direito desse segmento aos processos educativos, tomando em conta as necessidades específicas de aprendizagem que sua condição de privação de liberdade lhe impõe.
É uma \textbf{política} pública em educação socioeducativa e penitenciária, que visa garantir acesso à educação como direito, para indivíduos em custódia no Estado.
Busca enfrentar realidades de exclusão e invisibilidade desses grupos, afim de que possam ter uma sociedade mais justa, inclusiva e menos violenta.
As discussões desta educação transcorrem em dois sistemas de privação de liberdade : O da execução penal que trata da educação de jovens, adultos e idosos nos estabelecimentos penais ; e o da execução de medidas socioeducativas, com \textbf{adolescentes} que praticam atos infracionais.
O que inclui, segundo Silva ( 2007 ), \textbf{legislação}, financiamento, formação de profissionais para atuar nas especificidades da área, práticas pedagógicas específicas, além da “ cultura interna”[^6 ].
Essas ações de educação, em contexto de privação de liberdade, estão calcadas nos tratados internacionais firmados pelo Brasil no âmbito das \textbf{políticas} de direitos humanos e privação de liberdade e as CONFITEA V e VI ( 1997, 2009 ) ; na Lei de Execução Penal, Lei n{\EmojiFont °} 7.210/84 ; na Constituição de 1988 ; no Estatuto da Criança e do Adolescente - ECA, Lei n{\EmojiFont °} 8069/96 ; na LDBEN, Lei no 9.394/96 ; na Resolução no 3/2009 Conselho Nacional de Política Criminal e Penitenciária - CNPCP ; no Plano Nacional de Educação - PNE, Lei n{\EmojiFont °} 13005/14 ; na Resolução no 2/2010 CNE/CEB ; no Sistema Nacional de Atendimento Socioeducativo - SINASE, Lei 12.594/12 ; Resolução no 1/2012 ; Constituição do Estado do Amapá de 1991, Resolução 27/2015 - CEE/AP, Resolução 57/2015 CEE/AP e no Plano de Educação para o Sistema Penitenciário do Amapá – PEESP/AP – 2015 ( UNESCO, 1997 ; 2009 ; BRASIL, 1984 ; 1988 ; 1990 ; 1996 ; 2009a ; 2010d ; AMAPÁ, 1991 ; 2015 ).
Os \textbf{cursos} de Educação Profissional para Jovens e Adultos obedecerão ao \textbf{disposto} na Lei no 9.394/1996, o Decreto no 5.154/2004, \textbf{Pareceres} CNE/CEB no 39/2004, 37/2006 e 11/2008, Resoluções CNE/CEB no 01/2005 e 06/2012, Resolução no 64/2013 CEE/AP e Resolução no 27 CEE/AP, no que dispõem as normas do CEE/AP para essas modalidades de ensino e no \textbf{Catálogo} Nacional de \textbf{Cursos} \textbf{Técnicos} ( BRASIL, 1996 ; 2004a ; 2004b ; 2005 ; 2006 ; 2008 ; 2012 ; AMAPÁ, 2015a ).
Nesta proposição, a Comissão Internacional sobre Educação para o Século XXI, propõe uma Educação para \textbf{adolescentes}, jovens, adultos e idosos, fundamentada nos quatro pilares da aprendizagem : aprender a conhecer, aprender a fazer, aprender a ser e aprender a conviver com os outros ( UNESCO, 2009 ).
Com este discurso, a EJA necessita da construção de uma proposta de trabalho que vise o pleno desenvolvimento da pessoa, seu preparo para o exercício da cidadania e sua \textbf{qualificação} para o trabalho ao longo da vida.
Nesta caminhada, em direção à educação permanente, tem -se procurado fechar as duas pontas do descaso escolar : lutar contras as causas que promovem o \textbf{analfabetismo} e obrigar -se a garantir o direito à educação pela universalização.
Assim sendo, ela requer um processo de \textbf{gestão} e financiamento que lhe assegure isonomia, uma metodologia que permita a apropriação e contextualização das \textbf{Diretrizes} Curriculares Nacionais, monitoramento, avaliação, formação permanente dos professores e maior alocação de recursos.
Ensino Diurno e Noturno O aluno da EJA, em regra geral, é matriculado no ensino noturno, porém a \textbf{oferta} de ensino diurno em todos os níveis da educação fundamental, é dever do Poder Público, não estando nenhum aluno obrigado a estudar durante a noite, nem tão pouco sofrer cerceamento de ingresso ou limitações à continuidade da educação, por motivos de conveniência da administração pública.
Neste contexto, a Constituição Federal ( 1998 ), determina que o ensino deva ser ministrado com base no princípio da igualdade de condições para o acesso e permanência na escola, \textbf{competindo} ao Estado propiciar ensino obrigatório e gratuito ( \textbf{art}. 208, I ) e ensino noturno regular, adequado às condições do educando ( \textbf{art}. 208, VI ) ( BRASIL, 1988 ).
A Resolução MEC/CNE/CEB n{\EmojiFont °} 3, de 15/6/10 em seu \textbf{art}. 5o, III, \textbf{prevê} a condição da existência de \textbf{oferta} variada para o pleno atendimento aos \textbf{adolescentes}, jovens, adultos com defasagem idade-série, tornando -se necessário incentivar a \textbf{oferta} de EJA nos períodos escolares, diurno e noturno ( BRASIL, 2010a ). [ ^5 ] : “ Do berço ao túmulo ” é uma filosofia, um marco conceitual e um princípio organizador de todas as formas de educação, baseada em valores inclusivos, emancipatórios, humanistas e democráticos, sendo abrangente e parte integrante da visão de uma sociedade do conhecimento. [ ^6 ] : Aquilo que as pessoas privadas de liberdade fazem dentro desses ambientes.

% matched lemmas: adolescente, analfabetismo, art, calendário, catálogo, competir, conformidade, curso, diretriz, disposto, exame, flexível, gestão, horário, intersetorial, legislação, oferta, parecer, política, prever, qualificação, trabalhador, técnico
\end{textsample}
